\documentclass[12pt]{article}
\usepackage[utf8]{inputenc}
\usepackage[T1]{fontenc}
\usepackage{hyperref}
\title{Southern Sama Language Resources for K-12 Curriculum Development}
\author{Compiled by Claude Sonnet 5.5 (claude-sonnet-5-5) for Andr\'as Kornai}
\date{}
\begin{document}
\maketitle

\section{Introduction}
Southern Sama (ISO 639-3 ssb; also called Sinama Tawi-Tawi or Southern Sinama) is an Austronesian language of the Sama--Bajaw group. It is spoken in southern Sulu and Tawi-Tawi, Philippines, and also in Malaysia. Listed dialects include Balimbing, Bongao, Languyan, Obian, Sibutu, Simunul, Sitangkai and Tandubas. One source gives about 200,000 speakers, while the task description gives about 410,000. The Glottolog page found in search carries the identifier sout2844, which differs from the Glottocode given in the task description (sout2918). Please verify which is correct.

\section{Primary Educational Resources}
No Southern Sama primers, textbooks, or mother-tongue (MTB-MLE) materials were found in search. A direct inquiry to the Department of Education or SIL Philippines would be needed.

\section{Linguistic and Reference Materials}
\begin{itemize}
\item Verbix language page for Sama, Southern: \url{https://docs.verbix.com/Languages/SamaSouthern}
\item Glottolog entry (Glottocode as shown there; see note above): \url{https://glottolog.org/resource/languoid/id/sout2844}
\item sinama.org hosts a dictionary app, an orthography section and an educational poster. The dictionary described there is for Central Sinama, not Southern Sama, so use it for comparison only: \url{https://sinama.org/page/3/}
\item Overview of the Sama--Bajaw languages: \url{https://www.Wikipedia.org/wiki/Sinama}
\end{itemize}

\section{Religious and Community Resources}
Scripture Earth lists Bible materials for ssb, including an app for Sama Tawi-Tawi: \url{https://www.scriptureearth.org/00eng.php?iso=ssb}

\section{Digital and Audio Resources}
No dedicated audio or video collections were identified in search beyond the scripture resources above.

\section{Curriculum Development Context}
The task description says the language is used in education, but search found no confirming documents or materials. Related Sinama orthography and dictionary work for Central Sinama could serve as models, with dialect differences kept in mind.

\end{document}